\scanpage{18}

% 1–4, equivalence at head
\diagram{\begin{tikzcd}[ampersand replacement=\&,column sep=large]
\left.\begin{gathered}
 \text{géom. de drapeaux } D\\[-2pt]
 {}+\text{ ``\deleted{\unsure{application}}\inserted{fonction} dimension''}\\
 \delta:\underset{\substack{\downarrow\\\text{él. min.}}}{\Phi}\longrightarrow I
\end{gathered}\right| \arrow[r,Leftrightarrow] \&
\left|\ \text{foncteur }D'/\underline{\unsure{\mathscr P_f^*|\mathcal P_f^*|\mathfrak P_f^*}}(I)^o\to(\mathrm{Ens})\right.
\end{tikzcd}}

% 3–6, below right-hand member and across middle
\margin{\uline{NB} Posant $I_0=\{i\in I\mid D'(\{i\})\ne\varnothing\}$ on a sur $I_0$ une structure de \unsure{maquette simpl.|m. géométr. simpl.|ens. quasi-simpl.}, ayant comme simplexes les $\sigma$ tels que $D'(\sigma)\ne\varnothing$. \inserted{(NB $I_0=\delta(\Phi)$).}}

% 7–9, case of the natural numbers
Le cas qui nous intéresse plus particulièrement est le cas
\[
 I=\mathbb N
\]

% 10–14, dimension
auquel cas $\delta$ prend le nom de fonction dimension \deleted{\unsure{$\Phi/I$}} Et $\sup_{\overline{\mathbb R}}\delta(I)$ \deleted{\unsure{prend le nom}}\inserted{s'appelle} la dimension de la géom. de drapeaux $D$ \deleted{\unsure{avec sa dimension}}\inserted{(munie de sa} fonction dimension). Si celle-ci est finie, soit $n$, il y a souvent avantage à remplacer $\mathbb N$ par $[0,n]_{\mathbb N}$.

% 15–20, general construction
Plus gén., \inserted{(pour une géom. de drapeaux $D$ donnée)} supposons donnée une loi qui associe à chaque \deleted{partie} simplexe $\sigma\subset\underline{\Phi}$ (i.e. $\sigma\in\mathscr P_f^*(\Phi)$ de la forme $I_d$) un \uline{ordre total} sur $\sigma$, de façon que ces ordres totaux s'induisent mutuellement

% 20–27, orders and preorders
[ce qui est le cas pour les ordres totaux induits par un $\delta:\underline{\Phi}\to I$ si $I$ lui-m\unsure{e} est tot. ordonné, ou muni d'un \inserted{pré}ordre \deleted{\unsure{induisant}} pour lequel les simplexes de $I$ sont tot. ord...) mais cela revient à prendre sur $I$ le \inserted{pré}ordre image inverse...] Alors, posant pour $n\in\mathbb N$ ($n\geq0$)

% 28–30, D[n]
\[
 \begin{aligned}
 D[n]&=\{d\in D\mid\operatorname{card}I_d=n+1\}\\
 &=\coprod_{\sigma\in\mathscr P_{\unsure{f\,n+1|n+1}}(\Phi)}D(\sigma)\\
 &\left(\underset{\text{ou \unsure{s'ident.|seulement|aussi}}}{=}\
 \coprod_{\sigma\in\mathscr P_{f\,n+1}(I)}D'(\sigma)\right)
 \end{aligned}
\]

% lower left
% archive sheet number: 17
